\section{Background}
\label{sec:background}

\subsection{Infrasound and Human Physiology}

Infrasound refers to sound waves with frequencies below the nominal human
hearing threshold of \SI{20}{\hertz}. While inaudible in the conventional
sense, infrasound can be perceived through vibrotactile and whole-body
sensations at sufficiently high amplitudes \citep{Moller2004}.

% TODO: Review of infrasound perception thresholds
% TODO: Historical claims about the brown note

\subsection{Abdominal Cavity Mechanics}

The human abdominal cavity can be idealised as a fluid-filled elastic shell.
The abdominal wall consists of multiple tissue layers---skin, subcutaneous
fat, muscle (rectus abdominis, obliques, transversus), and fascia---that
together form a compliant enclosure \citep{Fung1993}.

% TODO: Typical material properties with ranges
% TODO: Geometric dimensions from anthropometric data

\subsection{Shell Vibration Theory}

For thin elastic shells, the natural frequencies depend on geometry, material
properties, and boundary conditions. The Donnell thin-shell equations provide
a classical framework for cylindrical shells \citep{Leissa1973, Soedel2004}.

% TODO: Key equations for cylindrical shell modes
% TODO: Extension to ellipsoidal geometries

\subsection{Acoustic Cavity Modes}

The natural frequencies of a rigid-walled acoustic cavity are governed by the
Helmholtz equation:
%
\begin{equation}
  \nabla^2 p + k^2 p = 0,
  \label{eq:helmholtz}
\end{equation}
%
where $p$ is the acoustic pressure and $k = \omega / c$ is the wavenumber.
For canonical geometries (rectangular, cylindrical, spherical), closed-form
solutions exist in terms of trigonometric and Bessel functions.

% TODO: Solutions for cylindrical and ellipsoidal cavities

\subsection{Fluid--Structure Interaction}

When the cavity walls are elastic rather than rigid, the structural vibration
and acoustic fields are coupled. The coupled eigenvalue problem takes the form:
%
\begin{equation}
  \begin{bmatrix}
    \mathbf{K}_s & \mathbf{C}_{fs} \\
    \mathbf{C}_{sf} & \mathbf{K}_f
  \end{bmatrix}
  \begin{Bmatrix}
    \mathbf{u} \\ \mathbf{p}
  \end{Bmatrix}
  = \omega^2
  \begin{bmatrix}
    \mathbf{M}_s & \mathbf{0} \\
    \mathbf{0} & \mathbf{M}_f
  \end{bmatrix}
  \begin{Bmatrix}
    \mathbf{u} \\ \mathbf{p}
  \end{Bmatrix},
  \label{eq:coupled_eigenvalue}
\end{equation}
%
where subscripts $s$ and $f$ denote the structural and fluid domains,
respectively, and $\mathbf{C}$ represents the coupling matrices
\citep{Junger1986}.
